\documentclass[10pt]{article}
\usepackage[margin=20mm]{geometry}
\usepackage{amsmath,amssymb,mathtools,fontspec,unicode-math,xurl,hyperref}
\IfFontExistsTF{DejaVu Serif}{\setmainfont{DejaVu Serif}}{\setmainfont{Latin Modern Roman}}
\setmathfont{latinmodern-math.otf}
\setlength{\parindent}{0pt}\setlength{\parskip}{6pt}\sloppy
\title{Cusp labels, height increments, and coding in diagonal-flow entropy bounds}
\author{AI-assisted technical review}
\date{}
\begin{document}
\maketitle
Abstract: The retained 29 December 2023 version of Ron Mor’s linear-functional method for diagonal-flow entropy is examined. The main substantive difficulty is an incompatibility between the printed reductive Borel–Serre convergence convention and the labeling of cusp regions by dimensions of short subspaces. A three-dimensional diagonal lattice contradicts two boundary identifications used in passing from finite-scale bounds to boundary masses. Locally repairable issues concern a uniform time interval, the sign of a height increment, a window-dependent coding identity, and two adapted-basis estimates. The finite-scale covering strategy is not refuted by these examples, and no counterexample involving invariant-measure entropy is asserted. Scientific adjudication remains pending.\par
\section*{Article, version, and mathematical scope}
The article is Ron Mor, Bounding entropy for one-parameter diagonal flows on \(\mathrm{SL}_d(\mathbb R)/\mathrm{SL}_d(\mathbb Z)\) using linear functionals, DOI 10.4171/JEMS/1592 [paper, journal]. The actual input is arXiv:2302.07122v2, dated 29 December 2023, with 55 PDF pages including references. The publisher records acceptance on 16 April 2024 and online publication on 13 February 2025. The restricted publisher PDF was not byte-compared with this author manuscript; equivalence is not assumed. All locations below refer to the retained PDF.\par
All 3,083 extracted-text lines and all 55 source pages were read or actually viewed, respectively. The 256-line manifest was read in full and its 13 listed inputs were checked by actual bytes and SHA256. Static source material was consulted in the bounded auxiliary checks. Source-archive possession does not establish scientific review of every archive member. No author code was executed, and no separate computational supplement was verified. Missing code is not a defect here.\par
Write \(G=\mathrm{SL}_d(\mathbb R)\), \(\Gamma=\mathrm{SL}_d(\mathbb Z)\), \(a_t=\exp(t\alpha)\), \(\sum_i\alpha_i=0\), and \(a=a_1\). For a standard upper block parabolic \(P\), the score is \[\chi_P(w)=h(P,a^w)-\phi\bigl(\pi_P(\alpha^w)\bigr),
 \qquad
 h(P,a^w)=\sum_{(i,j)\in E_P}
                   (\alpha_{w(i)}-\alpha_{w(j)})^+.\] Here \(\pi_P\) averages within blocks and the orientation is the double coset in \(\operatorname{stab}_W(a)\backslash W/W(T,P)\). The multiplicities of eigenvalues are not discarded.\par
Theorem 7.1 is the finite-scale entropy estimate. Theorems 1.1–1.2 translate it into a partition subordinate to boundary neighborhoods and a bound in terms of limiting boundary masses. Corollaries 1.3–1.4 concern total or stratum-specific escape of mass. Theorem 1.5 is a separate, explicit finite-scale consequence with \(\phi=0\). The optimized formulas and sharpness statements in Theorem 1.6 are attributed to subsequent work; they are not proved in this paper. Their external proofs have not been certified here.\par
\section*{Proof structure and affirmative checks}
The proof first controls diagonal trajectories of subspaces by grouping equal exterior-power weights. Its use of the Plücker relations addresses the fact that a dominant weight component need not itself be decomposable. Short rational subspaces are then isolated by gaps \(\eta_i=\lambda_i/\lambda_{i+1}\) in successive minima. The coding construction trims excursions at several scales, increasing the number of retained flag dimensions until it terminates after at most \(d\) steps.\par
Two estimates make the functional correction possible: the number of pieces is bounded by an arbitrarily small multiple of the observation time at suitably chosen cusp scales, and the sum of endpoint projection errors is \(O_a(rN)\). After discretization, the norm of the total \(\pi_P(\alpha^w)\) contribution is bounded by this error plus the difference between the endpoint heights. Adapted bases and parabolic Bowen-ball volume estimates supply a covering exponent. Birkhoff and Egorov then replace visiting frequencies by masses, with both endpoints in a large compact set. The entropy is affine under the discrete ergodic decomposition.\par
Several potential objections were explicitly checked and are not findings. The proof of Proposition 3.17 averages integer iterates and extends to an \(a\)-invariant measure, despite its stronger printed flow-invariance hypothesis. The subsequent Birkhoff and ergodic-decomposition arguments also concern \(T=a\). Lemma 7.3 is quoted with a stronger hypothesis and an omitted external proof; a matching discrete-time statement should be made explicit, but that wording difference alone is not a proved counterexample.\par
Negative individual values of \(\chi\) do not invalidate the visiting-frequency argument. Subtracting the global maximum leaves nonpositive coefficients, so frequency lower bounds give the required upper bound. Moreover the \(G\) score is \(h(G,a)\ge0\), as required in Proposition 4.9. For \(d=2\), \(\alpha=(1/2,-1/2)\) and \(\phi(h_1,h_2)=(h_1-h_2)/2\), both Borel orientation scores are exactly \(1/2\). The identity flow gives zero scores. These checks support the intended mechanism without resolving the boundary-identification problem below.\par
\section*{Substantive boundary-identification problem}
Problem 1 (Major: incompatible parabolic labels).  The convergence convention in §8.2, pp. 50–51, is incompatible with the short-gap labels used in the closure assertion on p. 52 and the set-limsup assertion on p. 54.\par
Decisive example. Let \(d=3\), let \(P_1=P_{(1,2)}\) and \(P_2=P_{(2,1)}\) be standard upper block parabolics, and set \[g_t=\operatorname{diag}(e^{2t},e^{-t},e^{-t}),\qquad x_t=g_t\mathbb Z^3.\] The printed convergence rule uses \(K M_P A_P N_P\) and requires \(\beta(a_j)\to\infty\) for every character in the adjoint action of \(A_P\) on \(\operatorname{Lie}N_P\). For \(P_1\), choose \(m_t=n_t=I\) and \(a_t=g_t\) in that decomposition. Both relevant characters equal \(e^{3t}\). Thus the printed rule gives \[x_t\longrightarrow y_\infty\in e_\infty(P_1),\] with fixed unit Levi data.\par
On the other hand, orthogonality of the rectangular lattice gives the exact successive minima \[(\lambda_1,\lambda_2,\lambda_3)
       =(e^{-t},e^{-t},e^{2t}),\qquad
 (\eta_1,\eta_2)=(1,e^{-3t}).\] Consequently, for any fixed sufficiently small \(0<\delta<\delta'\), eventually \[x_t\in\mathcal N_{\delta,\delta'}(P_2,P_2)\subset Y_{P_2}.\] Neither \(P_1\subset P_2\) nor \(P_2\subset P_1\) holds. They are not \(\Gamma\)-conjugate, since their stabilized flag dimensions differ. Therefore \[y_\infty\in\overline{Y_{P_2}}\cap e_\infty(P_1)\] contradicts the p. 52 closure assertion. For any admissible sequence of scales tending to zero, the same fixed \(y_\infty\) belongs to every \(\overline{\mathcal N_{\delta_j,\delta'_j}(P_2,P_2)}\). Hence the p. 54 identification of their set-limsup with \(e_\infty(P_2)\) also fails under the printed convention.\par
Impact and repair boundary. The latter identification is the stated bridge from (8.3) to the boundary-weighted conclusion of Theorem 1.2. It is a substantive proof issue even if the final theorem admits a repair. It does not supply an invariant probability measure, a measure-theoretic entropy calculation, or a counterexample to the numerical inequality. Likewise, accumulation at a boundary point by itself does not rule out an open set containing \(Y_{P_2}\) and omitting that point; the example directly refutes the printed closed-set assertions, not the bare open-cover existence statement of Theorem 1.1.\par
A correction must reconcile the left/right quotient and opposite-parabolic conventions with the short-subspace flag. If the printed convergence convention is retained, the relevant short dimensions are reversed by \(i\mapsto d-i\), corresponding to reversed block sizes. If the convention is changed, the boundary action and quotient topology must also be checked. In either case the authors should prove the corrected closure limits and verify which \(\chi_P\) coefficient multiplies each boundary mass. Merely renaming a symbol without these checks does not finish the repair. The finite-scale statement in Theorem 7.1 and the \(\phi=0\) consequence in Theorem 1.5 do not use this boundary identification.\par
\section*{Locally repairable findings}
Problem 2 (Minor: uniform time interval at the endpoint).  Proposition 3.13, p. 16, quantifies over all \(0<\delta<\eta_0\) while claiming a uniform time interval \(|t|\ll_a|\log\delta|\) on which \(\eta_l(a_t x)<\eta_0\).\par
Take \(d=2\), \(a_t=\operatorname{diag}(e^t,e^{-t})\), and \(x=\operatorname{diag}(s,s^{-1})\mathbb Z^2\), with \[s^2=\eta_0e^{-2\varepsilon},\qquad
 \delta=\eta_0e^{-\varepsilon},\qquad t=2\varepsilon.\] For sufficiently small positive \(\varepsilon\), the short direction does not interchange, and \[\eta_1(x)=\eta_0e^{-2\varepsilon}<\delta,\qquad
 \eta_1(a_t x)=\eta_0e^{2\varepsilon}>\eta_0,
 \qquad
 \frac{t}{|\log\delta|}\longrightarrow0.\] This disproves the uniform endpoint quantifier. Restricting to sufficiently small \(\delta\), for example \(\delta\le\eta_0^2\) with a suitable time constant, or using a time bound proportional to \(\log(\eta_0/\delta)\), repairs this obstruction. The elementary singular-value estimate \(\eta_l(a_t x)\le e^{2\|\alpha\|_\infty|t|}\eta_l(x)\) explains the latter scale. The later asymptotic use \(\delta\to0\) is not refuted.\par
Problem 3 (Minor: sign of the projected height increment).  The estimate immediately before (4.4) and (4.4) itself, p. 23, have the opposite sign from the definition \(\operatorname{height}(x)=\operatorname{diag}(-\log\lambda_i(x))\) on p. 21.\par
Proposition 3.18 gives \[\Delta\operatorname{height}(U,P)=-|U|\pi_P(\alpha^w)+O(1).\] Thus the bounded expression is the sum of these two terms, not their difference. For a direct test take the preceding flow, \(x=\operatorname{diag}(e^{-L},e^L)\mathbb Z^2\), \(P=B\), and \(U=[0,L/2]\). For large \(L\) this whole segment has a fixed short direction and remains in the required cusp. Exactly, \(\Delta\operatorname{height}(U,B)=-|U|\alpha\); the printed residual has norm \(2|U|\|\alpha\|\), which is unbounded. Changing the two signs repairs the formulas. Corollaries 4.5 and 4.7 use the norm of the total sum, so their intended cancellation estimate survives the sign correction. No change to the sign of \(\phi\) in the final entropy score is forced by this local correction.\par
Problem 4 (Minor: translating a coding is not rerunning it).  The equality \(a_{-N}Z(C,\delta,\delta',r)=Z_+(C',\delta,\delta',r)\) in the proof of Lemma 6.4, p. 44, fails because the excursion algorithm depends on the observation window.\par
Here is an exact log-coordinate example. Let \(k\) be a sufficiently large positive integer, and set \[\begin{aligned}
 a_t&=\operatorname{diag}(e^{t/100},e^{-t/100}),&
 x&=\operatorname{diag}(e^{-474k},e^{474k})\mathbb Z^2,\\
 \delta&=e^{-1000k},&
 \delta'&=e^{-250k},\\
 r&=\frac{49}{50},&
 N&=1100k .\end{aligned}\] The parameters satisfy \(\delta'>\sqrt\delta\), \(r^3>1/4\), and \(N>|\log\delta|\); choosing \(k\) large also ensures \(\delta'<\eta_0\). On the original interval \([-N,N]\), \[\log\eta_1(a_t x)=-948k+t/50
                      \in[-970k,-926k].\] There is no hit of either \(\delta\) or \(\delta^r=e^{-980k}\) in this window. The two-dimensional algorithm therefore labels the whole window by \(G\).\par
Now let \(y=a_{-N}x\). Its enlarged window \([-2N,2N]\) corresponds to the old times \([-3300k,1100k]\), on which \[\log\eta_1\in[-1014k,-926k].\] This interval does hit \(\delta\), while it remains entirely inside one \(\delta'\) excursion. Both relevant excursion labels are now \(\{1\}\), and the new algorithm labels the entire enlarged window by \(B\). The prescribed translate \(C'\) of the old constant-\(G\) coding is still constant \(G\). Thus \(y\) belongs to the left-hand set of the asserted equality and not the right.\par
Definition 13 allows all label-valued functions, so \(C'\in\mathcal C_{2N}\) itself is not the problem. The covering conclusion can instead be proved by stating Lemma 6.3 for a given partition and its prescribed continuous-time cusp labels. Its induction only uses those properties. Translate the old partition and labels, apply that induction directly, and retain the original bounds on their number and exceptional durations. This is a local proof repair; the example does not disprove the covering exponent.\par
Problem 5 (Minor: two adapted-basis normalizations).  Equation (5.1), p. 38, and the fixed-radius step of Proposition 6.2, p. 42, require small corrections.\par
In (5.1), \(l\) is the block number, determined by \(\dim Z_{l-1}<j\le\dim Z_l\). The summation must begin at \(\dim Z_l+1\), not \(l+1\). For \(d=2\), \(P=G\), \(u=I\), \(l=1\) and \(j=2\), the printed sum includes the diagonal entry of \(u\) and gives \(2z_2\) instead of \(z_2\). The corrected block index gives the intended transverse-component estimate.\par
Proposition 5.3 bounds \(\|a_t^w u_2a_{-t}^w-I\|_F\) at the fixed \(\varepsilon_0\) scale. It does not make \(wu_2w^{-1}\) an element of a Bowen ball of radius \(O(\rho)\) as \(\rho\) tends to zero, as asserted on p. 42. For example \(u=I+cE_{21}\) and \(a_t=\operatorname{diag}(e^t,e^{-t})\) give transverse entry \(ce^{-2t}\) and supremum \(c\) on \(t\ge0\). Keep a fixed small radius \(C\varepsilon_0\), then cover that Bowen box by \(O_{\rho,\varepsilon_0}(1)\) smaller Bowen boxes. The coordinate contraction factors cancel in the volume ratio; the extra factor is independent of \(T\). This retains the exponential covering rate.\par
For completeness, the nonnegativity justification on p. 53 can be stated using the fact that averaging \(\pi_H(\alpha^w)\) over permutations gives zero. Choose an orientation with \(\phi(\pi_H(\alpha^w))\le0\) to obtain \(h-\phi\ge0\). The printed existence of an orientation with \(\phi\ge0\) is true but has the wrong sign for that justification.\par
\section*{Supporting checks}
The recommended initial assessment is major revision because the boundary-mass identification supporting Theorem 1.2 requires a coherent convention and a replacement argument. Four further findings above are local and have explicit repair directions. This assessment does not require the entire article to be false or irreparable. No invariant-measure entropy counterexample was obtained.\par
The full primary proof chain, definitions, quantifiers, repeated eigenvalues, zero flow, signs of the functional contribution, clipped-window behavior, and boundary limits were examined. Classical geometry-of-numbers results, the external entropy-covering lemma, and the full Borel–Ji construction were not independently reproved from their original books or papers. In particular, the boundary counterexample is an internal test of the conventions actually printed here. Theorem 1.6’s follow-up claims remain outside this manuscript’s proved contribution.\par
The root executor performed the final scientific synthesis, wrote this English TeX report, ran the approved compiler, and actually viewed every retained source page and every final report page. Their read ranges and any read-only locating processes are separately recorded; neither compiled nor inspected page images. The root’s own saved exact-check script verifies the lattice exponents, endpoint limit, height sign, coding-window arithmetic, basis index, fixed-radius calculation, and rank-one score. These checks support the displayed examples; they are not a machine proof of the article.\par
The exact primary PDF contains 620,620 bytes. Its SHA256 is\par
> 20984c32231f0579f42b0ed6448ec204156ab40a55d7e92ea134c9486da9dcf2\par
The companion records bind the source manifest, input bytes, TeX, actual compiler and dependencies, build logs, PDF, rendered pages, and actual page-inspection receipts. No author-source compilation or author code was executed. No successful prior build was needlessly repeated.\par
This is the initial scientific report for M25-0219. The machine-readable scientific decision is `pending`; `human\_reviews` is empty. No independent acceptance or publication is claimed. The dispatcher’s observation of Personal/Chat/GPT-6 checked/ Pro 5-of-5 is attributed to that dispatcher. The runtime self-description is recorded separately, with backend identity and model equivalence unverified.\par
9 R. Mor, Bounding entropy for one-parameter diagonal flows on \(\mathrm{SL}_d(\mathbb R)/\mathrm{SL}_d(\mathbb Z)\) using linear functionals, arXiv:2302.07122v2, 29 December 2023. https://arxiv.org/abs/2302.07122v2. European Mathematical Society Press, article record for DOI 10.4171/JEMS/1592, accessed during this review. https://ems.press/journals/jems/articles/14298574.\par
\end{document}
